\chapter{Fitness for Older Adults}

\section*{Definition and Overview}
Fitness is just as important in later life as at any age, and its benefits are profound: regular physical activity prevents falls, maintains independence, protects against heart disease, diabetes, and dementia, and improves mood and quality of life. The Australian guidelines recommend that older adults do at least 30 minutes of moderate activity on most days, and emphasise four components: aerobic activity, strength training, balance exercises, and flexibility. Many older adults believe exercise is no longer for them or worry about injury, but the opposite is true: inactivity is the greater risk, and appropriate activity is safe for nearly everyone. This chapter explains how older adults can be active safely and effectively.

\section*{Epidemiology}
Physical activity declines with age, and a large proportion of Australians over 65 are insufficiently active. Inactivity in later life contributes to falls, which are the leading cause of injury-related hospitalisation and death in older Australians, to muscle loss (sarcopenia), osteoporosis, heart disease, and cognitive decline. Yet the majority of older adults could benefit from increased activity, and even small improvements in strength and balance reduce falls significantly. Structured programs such as group exercise classes and falls-prevention programs are available through community and health services and are effective at every age.

\section*{Aetiology and Pathophysiology}
Activity addresses the specific health threats of ageing.
\begin{itemize}
    \item \textbf{Falls:} balance exercises and strength training reduce the risk of falls and fractures, the most serious threat to independence in older age
    \item \textbf{Sarcopenia:} muscle mass and strength decline from around age 30, accelerating in later life; resistance training rebuilds and maintains muscle
    \item \textbf{Osteoporosis:} weight-bearing and resistance activity slows bone loss and protects against fractures
    \item \textbf{Cardiovascular health:} aerobic activity lowers blood pressure, cholesterol, and the risk of heart attack and stroke
    \item \textbf{Cognition:} regular activity is associated with a lower risk of dementia and slower cognitive decline
    \item \textbf{Mental health:} activity reduces depression, anxiety, and social isolation
    \item \textbf{Metabolic health:} activity improves blood sugar control and reduces the risk of type 2 diabetes
    \item \textbf{Independence:} maintaining strength, balance, and stamina allows older adults to manage daily activities and live independently
    \item \textbf{The four components:} aerobic (heart and lungs), strength (muscles and bone), balance (preventing falls), and flexibility (joints and movement)
\end{itemize}

\section*{Clinical Features}
An effective older-adult activity program has recognisable features.
\begin{itemize}
    \item Aerobic activity: brisk walking, swimming, cycling, or dancing, in sessions of at least 10 minutes, building to 30 minutes on most days
    \item Strength training: twice weekly, with resistance bands, light weights, or bodyweight exercises such as sit-to-stand and wall push-ups
    \item Balance exercises: standing on one foot, heel-to-toe walking, and tai chi, done regularly to prevent falls
    \item Flexibility: gentle stretching of the major muscle groups after activity or on rest days
    \item Progression that is gradual, with a rate that suits the individual's current fitness
    \item Activity that is enjoyable and social, such as group classes, walking groups, or swimming
    \item Appropriate footwear and a safe environment
    \item A plan that accounts for medical conditions and medicines
\end{itemize}

\section*{Differential Diagnosis}
Activity plans must be tailored to individual health circumstances.
\begin{itemize}
    \item Heart disease, diabetes, and lung disease, which need medical review and appropriate intensity
    \item Joint problems such as arthritis, which benefit from low-impact activity such as swimming and cycling
    \item Osteoporosis, which favours safe weight-bearing activity and strength work under guidance
    \item Balance problems and a history of falls, which require balance-specific programs and falls risk assessment
    \item Neurological conditions such as Parkinson's disease, which benefit from specialised exercise programs
    \item Cognitive impairment, where supervision and simple, structured activity are needed
    \item Medicines and conditions that affect heart rate and blood pressure during exercise
\end{itemize}

\section*{When to Seek Medical Care}
Older adults with chronic conditions, unexplained symptoms, or concerns about safety should discuss their activity plan with a doctor. A falls risk assessment is appropriate for anyone who has fallen, feels unsteady, or is starting balance exercise.

\textbf{Call 000 (or local emergency services) for:}
\begin{itemize}
    \item Chest pain, severe breathlessness, or dizziness during activity
    \item A fall with head injury, severe pain, or inability to move
    \item Fainting or collapse
    \item Any emergency symptoms
\end{itemize}

\section*{Diagnosis and Investigations}
\begin{itemize}
    \item \textbf{Medical review:} assessment of heart risk, blood pressure, diabetes control, and any symptoms on exertion
    \item \textbf{Falls risk assessment:} review of balance, muscle strength, vision, medicines, and home hazards
    \item \textbf{Bone density testing:} for women and men at risk of osteoporosis, guiding the safe exercise plan
    \item \textbf{Referral:} to physiotherapy for an individualised program, or to community exercise classes designed for older adults
    \item \textbf{Ongoing review:} adjusting the plan as fitness improves or health changes
\end{itemize}

\section*{Management}
\begin{itemize}
    \item \textbf{Start with what is achievable:} even 10-minute sessions several times a day count and build fitness
    \item \textbf{Use structured programs:} falls-prevention classes, group exercise, and tai chi are effective and social
    \item \textbf{Include strength twice a week:} chair stands, calf raises, and resistance bands protect muscle and bone
    \item \textbf{Practise balance regularly:} brief daily balance exercises prevent falls
    \item \textbf{Walk safely:} well-fitting shoes, a safe route, and a walking companion where useful
    \item \textbf{Manage the environment:} clear home hazards, improve lighting, and use handrails
    \item \textbf{Stay hydrated and dress appropriately} for weather and activity
    \item \textbf{Adapt for pain:} low-impact options such as water exercise suit arthritic joints
\end{itemize}

\section*{Complications}
\begin{itemize}
    \item Falls and fractures, which are more likely without balance and strength training
    \item Muscle loss and frailty from inactivity
    \item Worsening of heart disease and diabetes from a sedentary life
    \item Injuries from starting too vigorously or without appropriate technique
    \item Deconditioning after illness or hospitalisation if activity is not resumed
    \item Social isolation when mobility declines
\end{itemize}

\section*{Prognosis and Outcomes}
Older adults who are active live longer, healthier, and more independent lives. Falls risk falls by a third or more with balance and strength programs, muscle and bone are maintained, and the risk of heart disease, diabetes, and dementia is reduced. People who begin activity later in life still gain substantial benefit, so it is never too late to start. The most successful older exercisers are those who build activity into their routine, enjoy the social side, and progress gradually, and community programs make this achievable for nearly everyone.

\section*{Prevention and Self-Care}
\begin{itemize}
    \item Aim for 30 minutes of moderate activity on most days, including strength and balance work
    \item Join a group exercise or falls-prevention class
    \item Practise balance exercises daily, such as standing on one foot near a support
    \item Do strength exercises twice a week with bands, weights, or bodyweight
    \item Wear supportive, well-fitting shoes
    \item Keep the home safe: clear cords and rugs, improve lighting, and install handrails
    \item See a doctor before starting vigorous activity, and after any fall
    \item Involve family and friends to make activity social and enjoyable
\end{itemize}

\section*{Sources and Further Reading}
The following reputable sources provide further information for healthcare students and patients seeking reliable, current advice about fitness for older adults.
\begin{itemize}
    \item Department of Health and Aged Care. \textit{Physical activity guidelines for older Australians.} https://www.health.gov.au/
    \item Healthdirect Australia. \textit{Exercise for older Australians.} https://www.healthdirect.gov.au/
    \item Stay On Your Feet. \textit{Falls prevention programs.} https://www.stayonyourfeet.com.au/
    \item Better Health Channel. \textit{Exercise for seniors.} https://www.betterhealth.vic.gov.au/
    \item NHS. \textit{Physical activity for older adults.} https://www.nhs.uk/live-well/exercise/
    \item World Health Organization. \textit{Physical activity for older adults.} https://www.who.int/
\end{itemize}
